\section*{الفصل \ref{ch:g12:brain-and-movement} --- من النية إلى الحركة}

\begin{solution}{exo:g12:brain-and-movement:1}
شريط فوق قمة كل نصف كرة، أمام الأخدود الفاصل بين الفصين الجبهي
والجداري. وخريطته تسند كل جزء من الشريط إلى جزء من الجانب المقابل من
الجسم، بمساحة تتناسب مع دقة حركات ذلك الجزء.
\end{solution}

\begin{solution}{exo:g12:brain-and-movement:2}
محاوير \omterm{def:g12:stretch-reflex:neuron}{عصبونات} \omterm{def:g12:brain-and-movement:cortex}{القشرة الحركية} تنزل عبر جذع الدماغ، حيث يتقاطع معظمها
إلى الجانب الآخر، ثم تنزل في النخاع الشوكي إلى مستوى الساق، حيث تنتهي
على \omterm{def:g12:stretch-reflex:neuron}{العصبونات} الحركية (أو على \omterm{def:g12:stretch-reflex:neuron}{عصبونات} بينية مجاورة) التي تمتد محاويرها
إلى عضلة الساق.
\end{solution}

\begin{solution}{exo:g12:brain-and-movement:3}
المساحة تتبع دقة الضبط لا الحجم: فاليد تصنع حركات دقيقة مستقلة كثيرة
تحتاج \omterm{def:g12:stretch-reflex:neuron}{عصبونات} قشرية كثيرة؛ والجذع يصنع قليلًا منها وخشنًا.
\end{solution}

\begin{solution}{exo:g12:brain-and-movement:4}
كل أمر إلى عضلة --- من القشرة، ومن الأقواس المنعكسة، ومن \omterm{def:g12:stretch-reflex:neuron}{العصبونات}
البينية --- لا بد أن يمر عبر عصبوناتها الحركية، وهي تجمع كل الدخلات
ولا تطلق إلا إذا بلغ المجموع العتبة؛ فما من طريق آخر إلى العضلة.
\end{solution}

\begin{solution}{exo:g12:brain-and-movement:5}
مساحة الأصابع الموسَّعة عند الموسيقيين (وعند متطوعين بعد أيام من
التمرين)؛ واستيلاء الجيران على مساحة طرف مبتور؛ والتعافي بعد سكتة
تأمر به مناطق حول الآفة.
\end{solution}

\begin{solution}{exo:g12:brain-and-movement:6}
نحو \qty{4}{mV}؛ وأربعة دخلات متتابعة سريعًا تبلغ عتبة \qty{15}{mV}؛
والدخل المثبِّط الذي يصل مع دخل منبِّه يلغيه فلا يكاد الكمون يتحرك.
\end{solution}

\begin{solution}{exo:g12:brain-and-movement:7}
الوجه والذراع في اليسار، والساق أقل: أي \omterm{def:g12:brain-and-movement:cortex}{القشرة الحركية} في نصف الكرة
الأيمن، في جزئها الأسفل والأوسط (مساحتا الوجه والذراع)، ومساحة الساق
عند القمة ممسوسة جزئيًا لا غير --- وذلك فوق التقاطع، إذ كان العجز في
جانب واحد والوجه معه.
\end{solution}

\begin{solution}{exo:g12:brain-and-movement:8}
بعد السكتة تكون \omterm{def:g12:stretch-reflex:reflex}{القوس المنعكسة} سليمة لكن أثر القشرة المثبِّط في
\omterm{def:g12:stretch-reflex:neuron}{العصبونات} الحركية مفقود، فلا يلقى دخل المغزل إخمادًا: أي ركلة مبالغ
فيها. أما العصب المقطوع فيكسر القوس نفسها: فلا تبلغ إشارة العضلة، ولا
\omterm{def:g12:stretch-reflex:reflex}{منعكس}.
\end{solution}

\begin{solution}{exo:g12:brain-and-movement:9}
تفقد الحركة الإرادية في الساقين وأسفل الجذع (فالسبيل مقطوع)؛ وتبقى
\omterm{def:g12:stretch-reflex:reflex}{منعكسات} الساقين، مبالغًا فيها، إذ تقع أقواسها دون الآفة؛ وألياف الوجه
تغادر السبيل في جذع الدماغ، فوق الآفة، فلا تمس.
\end{solution}

\begin{solution}{exo:g12:brain-and-movement:10}
القبض يرسل تنبيهًا عامًا من القشرة نزولًا في النخاع، ومنه شيء إلى
\omterm{def:g12:stretch-reflex:neuron}{العصبونات} الحركية للساق؛ فيصير كمونها أقرب إلى العتبة، وتبلغها دخلات
المغزل أيسر وفي \omterm{def:g12:stretch-reflex:neuron}{عصبونات} أكثر.
\end{solution}

\begin{solution}{exo:g12:brain-and-movement:11}
سنوات من التمرين وسّعت مساحة الأصابع باللدونة؛ والبدء قبل السابعة، في
الطور الذي تعيد فيه القشرة تنظيم نفسها أسهل ما يكون، أنتج تغيرًا أكبر
مما ينتجه التمرين نفسه لاحقًا.
\end{solution}

\begin{solution}{exo:g12:brain-and-movement:12}
الدارات تقوى بالاستعمال: فإذا فعلت الذراع السليمة كل شيء، لم تنشَّط
الدارات الناجية إلى الذراع المشلولة قط ولم تنتقَ؛ وإكراه استعمالها
يجعلها تطلق، وما ينجح منها يقوى. وكما في الهرة، الوصلات المستعملة تأخذ
الإقليم؛ وغير المستعملة تفقده.
\end{solution}

\begin{solution}{exo:g12:brain-and-movement:13}
المساحة الحسية لليد، وقد حرمت الدخل، غزتها مساحة الوجه المجاورة، فصار
لمس الوجنة ينشّط \omterm{def:g12:stretch-reflex:neuron}{عصبونات} ما زالت "تعني" اليد. والخريطة الحركية تقع إلى
جانبها وتعيد تنظيم نفسها على النحو نفسه: فيتوقع أن تستولي الذراع
والوجه على مساحة اليد السابقة.
\end{solution}

\begin{solution}{exo:g12:brain-and-movement:14}
تخريب \omterm{def:g12:stretch-reflex:neuron}{العصبونات} الحركية: شلل، \omterm{def:g12:stretch-reflex:reflex}{ومنعكسات} غائبة، وعضلات ضامرة (فلا دخل
عصبي)؛ والقشرة تكوّن نيات لا تستطيع أن تبلغ أي عضلة --- فواجهة تقرأ
القشرة تستطيع أن تقود طرفًا اصطناعيًا أو منبّهًا للعضلات. وتخريب القشرة:
شلل، \omterm{def:g12:stretch-reflex:reflex}{ومنعكسات} مبالغ فيها، وعضلات تبقيها \omterm{def:g12:stretch-reflex:reflex}{المنعكسات}؛ \omterm{def:g12:stretch-reflex:neuron}{والعصبونات} الحركية
سليمة --- فتستطيع واجهة أن تنبه النخاع أو \omterm{def:g12:stretch-reflex:neuron}{العصبونات} مباشرة، لكن النية
نفسها هي المفقودة، ولدونة القشرة المحيطة هي الأمل الأكبر.
\end{solution}

\begin{solution}{exo:g12:brain-and-movement:15}
العضلات تنمو \omterm{def:g10:sport-and-health:training}{بالتدريب}، لكن الحركة المتعلَّمة تغير في الجهاز العصبي:
فالمساحة القشرية التي تأمر بها تتسع ووصلاتها تشحذ (وهي اللدونة)، ونمط
الدخلات البالغة كل \omterm{def:g12:stretch-reflex:neuron}{عصبون} حركي --- أيها يطلق، وكم واحدًا، وبأي ترتيب ---
يهذَّب حتى تنتج مجاميع \omterm{def:g12:stretch-reflex:neuron}{العصبونات} الحركية التقلص الصحيح في اللحظة
الصحيحة. فالعضلات هي هي، والأمر أفضل.
\end{solution}

\begin{solution}{pb:g12:brain-and-movement:1}
\textbf{1.} الألياف القشرية النخاعية تتقاطع في جذع الدماغ: فالقشرة
اليسرى تأمر الجانب الأيمن.

\textbf{2.} الجزآن الأسفل والأعلى من الشريط كله: الوجه والذراع (جانب
الشريط) والساق (قمته) --- أي آفة تعبر \omterm{def:g12:brain-and-movement:cortex}{القشرة الحركية} اليسرى.

\textbf{3.} الشريط الحسي، خلف الحركي مباشرة عبر الأخدود المركزي، سليم:
فالإحساس يبلغه على نحو طبيعي.

\textbf{4.} \omterm{def:g12:stretch-reflex:reflex}{القوس المنعكسة} في النخاع سليمة، ودخلات القشرة المخمِّدة إلى
عصبوناتها الحركية زالت: فدخل المغزل يطلق \omterm{def:g12:stretch-reflex:neuron}{العصبونات} أيسر.

\textbf{5.} \omterm{def:g12:stretch-reflex:neuron}{العصبونات} الحركية أو أعصابها في اليمين: فالقوس مكسورة (فلا
\omterm{def:g12:stretch-reflex:reflex}{منعكس}) والعضلات، بلا دخل عصبي، تضمر. أي آفة دون \omterm{prop:g12:brain-and-movement:integration}{السبيل النهائي
المشترك} لا فوقه.

\textbf{6.} $12 \times 1 + 3 \times 2 = \qty{18}{mV}$: أي فوق العتبة،
فقد أطلق.

\textbf{7.} $18 - 4 = \qty{14}{mV}$: أي دون العتبة بقليل، فلا إطلاق
--- والتثبيط عدّل الخطوة.

\textbf{8.} $0 + 6 = \qty{6}{mV}$: فلا إطلاق؛ والساق لا تتحرك.

\textbf{9.} قبلها: $20 - 6 = \qty{14}{mV}$، أي دون العتبة أو عندها
بالضبط --- فركلة متواضعة. وبعدها: \qty{20}{mV}، أي فوقها بكثير ---
فركلة قوية: فالإخماد هو ما كان يبقي \omterm{def:g12:stretch-reflex:reflex}{المنعكس} متواضعًا.

\textbf{10.} $8 + 6 = \qty{14}{mV}$: أي أقل بقليل؛ فمع قليل زائد من
دخل المغزل (أو من الجهد) يطلق. والحركة ضعيفة مترددة، تطلق في بعض
المحاولات دون بعض.

\textbf{11.} من القشرة الناجية حول الآفة ومن نصف الكرة الآخر، وعصبوناته
تكوّن وصلات جديدة أو مقواة إلى \omterm{def:g12:stretch-reflex:neuron}{العصبونات} الحركية للساق.

\textbf{12.} اللدونة تقوّي الدارات التي تنشَّط: فالحركة المحاولة تطلق
\omterm{def:g12:stretch-reflex:neuron}{العصبونات} القشرية الناجية وسبلها، وما ينتج منها تقلصًا ينتقى؛ أما الساق
التي تحرَّك تحريكًا سلبيًا فتنشّط الدارات الحسية لا الآمرة.

\textbf{13.} الكسر الصغير من الألياف القشرية النخاعية التي لا تتقاطع،
والسبل المارة عبر جذع الدماغ، تستطيع أن تحمل أمرًا من نصف الكرة الأيمن
إلى الساق اليمنى؛ وإعادة التنظيم تستقدمها.

\textbf{14.} اليد تشغل مساحة كبيرة متخصصة يعسر على المناطق المجاورة أن
تعوّض ضبطها الدقيق، وحركاتها تقتضي دقة أكبر بكثير من خطوة؛ أما أوامر
الساق الأخشن فأيسر في إعادة البناء.

\textbf{15.} المشترك: أن التعافي يتوقف على الاستعمال، وأن الدارات غير
المستعملة تفقد. والمختلف: أن القشرة البالغة تعيد تنظيم نفسها أبطأ بكثير
وأنقص مما تفعل قشرة الهرة داخل طورها الحرج --- ولهذا أخذ تعافيها أشهرًا
وبقي جزئيًا.

\textbf{16.} \omterm{def:g12:brain-and-movement:cortex}{القشرة الحركية}، \omterm{prop:g12:brain-and-movement:pathway}{والسبيل القشري النخاعي}، والتقاطع في جذع
الدماغ، \omterm{def:g12:stretch-reflex:neuron}{والعصبون} الحركي في النخاع (وهو الملتقى مع دخل \omterm{def:g12:stretch-reflex:reflex}{المنعكس})،
والأستيل كولين عند الملتقى العصبي العضلي، \omterm{def:g12:respiration-fermentation:atp}{وATP} الآتي من متقدرات ألياف
العضلة.

\textbf{17.} \omterm{prop:g12:glucose-and-diabetes:hormones}{الغلوكاغون}؛ \omterm{prop:g12:glucose-and-diabetes:hormones}{وجزر البنكرياس}.

\textbf{18.} توقع الدماغ (فالأمر ينسخ إلى المراكز القلبية، ومعه
\omterm{prop:g10:heart-lungs-effort:control}{الأدرينالين}) \omterm{def:g11:hormones-and-reproduction:hormone}{والتغذية الراجعة} من مستقبلات العضلات والدم.

\textbf{19.} اللدونة: فصلا الرؤية والدماغ، وهذا الفصل؛ والشريان
المسدود: فصلا القلب \omterm{def:g10:sport-and-health:training}{والنشاط البدني} والصحة؛ والمناعة: فصلا \omterm{def:g12:innate-immunity:immunity}{المناعة
الفطرية} والتكيّفية.

\textbf{20.} \omterm{def:g12:brain-and-movement:cortex}{القشرة الحركية الأولية} في نصف الكرة الأيسر؛ ومجموع دخلات
يبلغ \qty{15}{mV} فوق الراحة في غضون ميلي ثوان قليلة؛ ولدونة القشرة.
\end{solution}
